% Editable TikZ: boundary values, image energy, polarization, and circuits.
% Deterministic editable paths; see make_electrostatic_tikz.py.
\newcommand{\ThUniformSphereAnalysisDrawing}{\begin{tikzpicture}[>=Latex,font=\small,x=1.15cm,y=1.15cm]
\begin{scope}\clip(-3.2,-2)rectangle(3.2,2);
\draw[TextGray,dashed,line width=.55pt] plot coordinates{(2.08642,-2.00000) (2.18511,-0.64000) (2.20557,0.00000) (2.18511,0.64000) (2.08642,2.00000)};
\draw[TextGray,dashed,line width=.55pt] plot coordinates{(1.39621,-2.00000) (1.45922,-1.50000) (1.62487,-0.53333) (1.65287,-0.26000) (1.66202,0.00000) (1.65287,0.26000) (1.62487,0.53333) (1.45922,1.50000) (1.39621,2.00000)};
\draw[TextGray,dashed,line width=.55pt] plot coordinates{(0.67151,-2.00000) (0.71902,-1.67333) (0.79019,-1.40000) (0.87304,-1.18667) (1.08800,-0.72311) (1.17671,-0.47333) (1.22840,-0.23333) (1.24508,0.00000) (1.22840,0.23333) (1.17671,0.47333) (1.08800,0.72311) (0.87304,1.18667) (0.79019,1.40000) (0.71902,1.67333) (0.67151,2.00000)};
\draw[TextGray,dashed,line width=.55pt] plot coordinates{(0.00000,-2.00000) (0.00000,-1.00667)};
\draw[TextGray,dashed,line width=.55pt] plot coordinates{(0.00000,1.00667) (0.00000,2.00000)};
\draw[TextGray,dashed,line width=.55pt] plot coordinates{(-0.67151,-2.00000) (-0.71902,-1.67333) (-0.79019,-1.40000) (-0.87304,-1.18667) (-1.08800,-0.72311) (-1.17671,-0.47333) (-1.22840,-0.23333) (-1.24508,0.00000) (-1.22840,0.23333) (-1.17671,0.47333) (-1.08800,0.72311) (-0.87304,1.18667) (-0.79019,1.40000) (-0.71902,1.67333) (-0.67151,2.00000)};
\draw[TextGray,dashed,line width=.55pt] plot coordinates{(-1.39621,-2.00000) (-1.45922,-1.50000) (-1.62487,-0.53333) (-1.65287,-0.26000) (-1.66202,0.00000) (-1.65287,0.26000) (-1.62487,0.53333) (-1.45922,1.50000) (-1.39621,2.00000)};
\draw[TextGray,dashed,line width=.55pt] plot coordinates{(-2.08642,-2.00000) (-2.18511,-0.64000) (-2.20557,0.00000) (-2.18511,0.64000) (-2.08642,2.00000)};
\draw[DeepBlue,line width=.7pt] plot coordinates{(3.20000,-0.33644) (2.36800,-0.32350) (1.79219,-0.30000) (1.35200,-0.26065) (0.98400,-0.20085)};
\draw[->,DeepBlue,line width=.7pt] (1.82812,-0.30259)--(2.02788,-0.31246);
\draw[DeepBlue,line width=.7pt] plot coordinates{(-0.98400,-0.20085) (-1.35200,-0.26065) (-1.79219,-0.30000) (-2.36800,-0.32350) (-3.20000,-0.33644)};
\draw[->,DeepBlue,line width=.7pt] (-2.25194,-0.32049)--(-2.05206,-0.31342);
\draw[DeepBlue,line width=.7pt] plot coordinates{(0.98400,0.20085) (1.35200,0.26065) (1.79219,0.30000) (2.36800,0.32350) (3.20000,0.33644)};
\draw[->,DeepBlue,line width=.7pt] (2.05206,0.31342)--(2.25194,0.32049);
\draw[DeepBlue,line width=.7pt] plot coordinates{(-3.20000,0.33644) (-2.36800,0.32350) (-1.79219,0.30000) (-1.35200,0.26065) (-0.98400,0.20085)};
\draw[->,DeepBlue,line width=.7pt] (-2.02788,0.31246)--(-1.82812,0.30259);
\draw[DeepBlue,line width=.7pt] plot coordinates{(3.20000,-0.58328) (2.33600,-0.56090) (1.75200,-0.52074) (1.31200,-0.45388) (1.12000,-0.40611) (0.94400,-0.34858)};
\draw[->,DeepBlue,line width=.7pt] (1.77139,-0.52340)--(1.97065,-0.54067);
\draw[DeepBlue,line width=.7pt] plot coordinates{(-0.94400,-0.34858) (-1.12000,-0.40611) (-1.31200,-0.45388) (-1.75200,-0.52074) (-2.33600,-0.56090) (-3.20000,-0.58328)};
\draw[->,DeepBlue,line width=.7pt] (-2.19581,-0.55458)--(-1.99619,-0.54236);
\draw[DeepBlue,line width=.7pt] plot coordinates{(0.94400,0.34858) (1.12000,0.40611) (1.31200,0.45388) (1.75200,0.52074) (2.33600,0.56090) (3.20000,0.58328)};
\draw[->,DeepBlue,line width=.7pt] (1.99619,0.54236)--(2.19581,0.55458);
\draw[DeepBlue,line width=.7pt] plot coordinates{(-3.20000,0.58328) (-2.33600,0.56090) (-1.75200,0.52074) (-1.31200,0.45388) (-1.12000,0.40611) (-0.94400,0.34858)};
\draw[->,DeepBlue,line width=.7pt] (-1.97065,0.54067)--(-1.77139,0.52340);
\draw[DeepBlue,line width=.7pt] plot coordinates{(3.20000,-0.87086) (2.28000,-0.83787) (1.94400,-0.81180) (1.66400,-0.77768) (1.42400,-0.73358) (1.21600,-0.67782) (1.03200,-0.60805) (0.86332,-0.52057)};
\draw[->,DeepBlue,line width=.7pt] (1.68078,-0.78117)--(1.87922,-0.80603);
\draw[DeepBlue,line width=.7pt] plot coordinates{(-0.86332,-0.52057) (-1.03200,-0.60805) (-1.21600,-0.67782) (-1.42400,-0.73358) (-1.66400,-0.77768) (-1.94400,-0.81180) (-2.28000,-0.83787) (-3.20000,-0.87086)};
\draw[->,DeepBlue,line width=.7pt] (-2.11561,-0.82726)--(-1.91639,-0.80968);
\draw[DeepBlue,line width=.7pt] plot coordinates{(0.86332,0.52057) (1.03200,0.60805) (1.21600,0.67782) (1.42400,0.73358) (1.66400,0.77768) (1.94400,0.81180) (2.28000,0.83787) (3.20000,0.87086)};
\draw[->,DeepBlue,line width=.7pt] (1.91639,0.80968)--(2.11561,0.82726);
\draw[DeepBlue,line width=.7pt] plot coordinates{(-3.20000,0.87086) (-2.28000,0.83787) (-1.94400,0.81180) (-1.66400,0.77768) (-1.42400,0.73358) (-1.21600,0.67782) (-1.03200,0.60805) (-0.86332,0.52057)};
\draw[->,DeepBlue,line width=.7pt] (-1.87922,0.80603)--(-1.68078,0.78117);
\draw[DeepBlue,line width=.7pt] plot coordinates{(3.20000,-1.15426) (2.64000,-1.13632) (2.19200,-1.11072) (1.83200,-1.07632) (1.52800,-1.03000) (1.27200,-0.96973) (1.06400,-0.89627) (0.88800,-0.80466) (0.73500,-0.68750)};
\draw[->,DeepBlue,line width=.7pt] (1.56119,-1.03732)--(1.75881,-1.06804);
\draw[DeepBlue,line width=.7pt] plot coordinates{(-0.73500,-0.68750) (-0.88800,-0.80466) (-1.06400,-0.89627) (-1.27200,-0.96973) (-1.52800,-1.03000) (-1.83200,-1.07632) (-2.19200,-1.11072) (-2.64000,-1.13632) (-3.20000,-1.15426)};
\draw[->,DeepBlue,line width=.7pt] (-2.01542,-1.09662)--(-1.81658,-1.07519);
\draw[DeepBlue,line width=.7pt] plot coordinates{(0.73500,0.68750) (0.88800,0.80466) (1.06400,0.89627) (1.27200,0.96973) (1.52800,1.03000) (1.83200,1.07632) (2.19200,1.11072) (2.64000,1.13632) (3.20000,1.15426)};
\draw[->,DeepBlue,line width=.7pt] (1.81658,1.07519)--(2.01542,1.09662);
\draw[DeepBlue,line width=.7pt] plot coordinates{(-3.20000,1.15426) (-2.64000,1.13632) (-2.19200,1.11072) (-1.83200,1.07632) (-1.52800,1.03000) (-1.27200,0.96973) (-1.06400,0.89627) (-0.88800,0.80466) (-0.73500,0.68750)};
\draw[->,DeepBlue,line width=.7pt] (-1.75881,1.06804)--(-1.56119,1.03732);
\draw[DeepBlue,line width=.7pt] plot coordinates{(3.20000,-1.38195) (2.58400,-1.36064) (2.10400,-1.33059) (1.70400,-1.28821) (1.37600,-1.23213) (1.10400,-1.15963) (0.87863,-1.06667) (0.70457,-0.95333) (0.57985,-0.82000)};
\draw[->,DeepBlue,line width=.7pt] (1.43760,-1.24593)--(1.63477,-1.27945);
\draw[DeepBlue,line width=.7pt] plot coordinates{(-0.57985,-0.82000) (-0.70457,-0.95333) (-0.87863,-1.06667) (-1.10400,-1.15963) (-1.37600,-1.23213) (-1.70400,-1.28821) (-2.10400,-1.33059) (-2.58400,-1.36064) (-3.20000,-1.38195)};
\draw[->,DeepBlue,line width=.7pt] (-1.90952,-1.31333)--(-1.71089,-1.28992);
\draw[DeepBlue,line width=.7pt] plot coordinates{(0.57985,0.82000) (0.70457,0.95333) (0.87863,1.06667) (1.10400,1.15963) (1.37600,1.23213) (1.70400,1.28821) (2.10400,1.33059) (2.58400,1.36064) (3.20000,1.38195)};
\draw[->,DeepBlue,line width=.7pt] (1.71089,1.28992)--(1.90952,1.31333);
\draw[DeepBlue,line width=.7pt] plot coordinates{(-3.20000,1.38195) (-2.58400,1.36064) (-2.10400,1.33059) (-1.70400,1.28821) (-1.37600,1.23213) (-1.10400,1.15963) (-0.87863,1.06667) (-0.70457,0.95333) (-0.57985,0.82000)};
\draw[->,DeepBlue,line width=.7pt] (-1.63477,1.27945)--(-1.43760,1.24593);
\draw[DeepBlue,line width=.7pt] plot coordinates{(3.20000,-1.63843) (2.50400,-1.61339) (1.96000,-1.57719) (1.50400,-1.52532) (1.12800,-1.45649) (0.79200,-1.36056) (0.64800,-1.30315) (0.52800,-1.24319) (0.42842,-1.18000) (0.34924,-1.11333) (0.29027,-1.04000) (0.25938,-0.97052)};
\draw[->,DeepBlue,line width=.7pt] (1.21760,-1.47690)--(1.41440,-1.51256);
\draw[DeepBlue,line width=.7pt] plot coordinates{(-0.25938,-0.97052) (-0.29027,-1.04000) (-0.34924,-1.11333) (-0.42842,-1.18000) (-0.52800,-1.24319) (-0.64800,-1.30315) (-0.79200,-1.36056) (-1.12800,-1.45649) (-1.50400,-1.52532) (-1.96000,-1.57719) (-2.50400,-1.61339) (-3.20000,-1.63843)};
\draw[->,DeepBlue,line width=.7pt] (-1.72482,-1.55438)--(-1.52640,-1.52926);
\draw[DeepBlue,line width=.7pt] plot coordinates{(0.25938,0.97052) (0.29027,1.04000) (0.34924,1.11333) (0.42842,1.18000) (0.52800,1.24319) (0.64800,1.30315) (0.79200,1.36056) (1.12800,1.45649) (1.50400,1.52532) (1.96000,1.57719) (2.50400,1.61339) (3.20000,1.63843)};
\draw[->,DeepBlue,line width=.7pt] (1.52640,1.52926)--(1.72482,1.55438);
\draw[DeepBlue,line width=.7pt] plot coordinates{(-3.20000,1.63843) (-2.50400,1.61339) (-1.96000,1.57719) (-1.50400,1.52532) (-1.12800,1.45649) (-0.79200,1.36056) (-0.64800,1.30315) (-0.52800,1.24319) (-0.42842,1.18000) (-0.34924,1.11333) (-0.29027,1.04000) (-0.25938,0.97052)};
\draw[->,DeepBlue,line width=.7pt] (-1.41440,1.51256)--(-1.21760,1.47690);
\draw[DeepBlue,line width=.7pt] plot coordinates{(3.20000,-1.96324) (2.36000,-1.93307) (1.68800,-1.88711) (1.12800,-1.82356) (0.31200,-1.69663) (0.00000,-1.67513) (-0.31200,-1.69663) (-1.12800,-1.82356) (-1.68800,-1.88711) (-2.36000,-1.93307) (-3.20000,-1.96324)};
\draw[->,DeepBlue,line width=.7pt] (-0.42417,-1.71118)--(-0.22583,-1.68544);
\draw[DeepBlue,line width=.7pt] plot coordinates{(-3.20000,1.96324) (-2.36000,1.93307) (-1.68800,1.88711) (-1.12800,1.82356) (-0.31200,1.69663) (0.00000,1.67513) (0.31200,1.69663) (1.12800,1.82356) (1.68800,1.88711) (2.36000,1.93307) (3.20000,1.96324)};
\draw[->,DeepBlue,line width=.7pt] (0.22583,1.68544)--(0.42417,1.71118);
\end{scope}
\fill[DeepBlue!6](0,0)circle(1);\draw[thick](0,0)circle(1);\node at(0,0){$E=0$};
\foreach\ang in {-60,-30,0,30,60}{\node at({.88*cos(\ang)},{.88*sin(\ang)}){$+$};\node at({-.88*cos(\ang)},{.88*sin(\ang)}){$-$};}
\draw[->,DeepBlue,thick](-3,2.3)--(-1.9,2.3)node[midway,above]{$\vect E_0$};\draw[->](1.8,-2.35)--(3,-2.35)node[right]{$+z$};
\node[below,align=center]at(-.3,-2.25){$r=R$：$\varphi=0,\ E_\theta=0$\\$\sigma=3\varepsilon_0E_0\cos\theta$};
\node[above]at(.7,2.15){蓝实线为场线；灰虚线为等势线};
\end{tikzpicture}}

\providecommand{\ThermalFigure}[3]{\begin{qualityfigure}
\begin{tikzpicture}[>=Latex,font=\small,x=.8cm,y=.8cm]#1\end{tikzpicture}
\figcap{#2}\label{#3}\end{qualityfigure}}
\newcommand{\ThConductorsBase}[2]{
 \draw[fill=DeepBlue!8,thick](-2,0)ellipse(1.3 and 1);\draw[fill=DeepBlue!8,thick](2,0)ellipse(.9 and 1.2);
 \node at(-2,0){$A:#1$};\node at(2,0){$B:#2$};\node[above]at(0,1.9){形状、尺寸和间距不预设对称};}
\newcommand{\ThQOppositeConductors}{\ThermalFigure{\ThConductorsBase{+Q}{-q}}{两个孤立导体的给定净电荷，不预先指定局部表面电荷的符号。}{ch20:quality:opposite-q}}
\newcommand{\ThQPositiveConductors}{\ThermalFigure{\ThConductorsBase{+Q}{+q}}{两导体净电荷均为正，电势大小关系还依赖几何。}{ch20:quality:positive-q}}
\newcommand{\ThAPositiveConductors}{\ThermalFigure{
 \draw[->](0,0)--(7,0)node[right]{从高电势导体表面向外};\draw[->](0,0)--(0,3)node[above]{$\varphi$};
 \draw[DeepBlue,thick](0,2.5)..controls(2,1.8)and(3,.6)..(6.5,.2);\node[left]at(0,2.5){$V_{\max}$};
 \node[right]at(3.5,2.2){$E_n=-\partial_n\varphi\ge0$};
}{最高电势边界的向外导数非正；仅给净电荷大小不能先指定谁处于最高电势。}{ch20:quality:positive-a}
图~\ref{ch20:quality:positive-a}从最高电势导体的表面沿外法线出发。拉普拉斯方程的最大值原理禁止外部电势超过该最大边界值，故其面电荷非负；较低电势导体的局部导数可以反号，不能由净电荷正号直接确定局部符号。
}
\newcommand{\ThQNeutralConductors}{\ThermalFigure{\ThConductorsBase{0}{0}
 \node[below]at(0,-1.7){无外加电荷、无外场，两球初始各自中性};
}{考察静电平衡能否仅由相互感应自发产生持续带电分布。}{ch20:quality:neutral-q}}
\newcommand{\ThANeutralConductors}{\ThermalFigure{
 \draw[thick](-2,0)circle(1);\draw[thick](2,0)circle(1);\draw[dashed](-4,-1.8)rectangle(4,1.8);
 \node at(-2,0){$Q_A=0$};\node at(2,0){$Q_B=0$};\node[above]at(0,2.1){$\varphi(\infty)=0$};
 \node[below]at(0,-2.1){候选平衡解：$\varphi=0$，$\vect E=0$};
}{唯一性问题的边界条件：各导体等势且各自净电荷为零。}{ch20:quality:neutral-a}
图~\ref{ch20:quality:neutral-a}中给定的是两个独立的净电荷约束，不是允许两球交换任意电荷的单一约束。处处零电势同时满足全部条件；唯一性定理排除另外一组稳定、非零的经典静电场。
}
\newcommand{\ThQLikeAttraction}{\ThermalFigure{
 \draw[fill=DeepBlue!10,thick](0,0)circle(1.4);\node at(0,0){净电荷 $+q$};\fill[DeepBlue](4,0)circle(3pt)node[above]{$+q$};
 \draw[<->](0,-2)--node[below]{$a$}(4,-2);\draw[->](0,0)--(120:1.4)node[midway,left]{$R$};
}{真实系统只有一个孤立导体球和一个球外点电荷，$a>R$。}{ch20:quality:like-q}}
\newcommand{\ThQImageEnergy}{\ThermalFigure{
 \draw[thick](0,0)circle(1.2);\fill[DeepBlue](3,0)circle(3pt)node[above]{$+Q$};\draw[<->](0,-1.7)--node[below]{$L>R$}(3,-1.7);
 \node[below]at(0,-2.4){接地球外};\begin{scope}[shift={(7,0)}]\draw[thick](0,0)circle(1.7);\draw[thick](0,0)circle(1.2);\fill(.6,0)circle(3pt)node[above]{$-q$};\draw[<->](0,-.5)--node[below]{$d<R$}(.6,-.5);\node[below]at(0,-2.4){接地球腔内};\end{scope}
}{比较外部与内部点电荷的有限感应能；点电荷自身发散的自能须扣除。}{ch20:quality:image-energy-q}}
\newcommand{\ThAImageEnergy}{\ThermalFigure{
 \draw[->](0,0)--(6,0)node[right]{建立电荷的比例 $\lambda$};\draw[->](0,0)--(0,3)node[above]{$|q\varphi_{\rm ind}(\lambda)|$};
 \draw[DeepBlue,thick](0,0)--(5,2.5);\fill[DeepBlue!12](0,0)--(5,0)--(5,2.5)--cycle;\node[below]at(5,0){$1$};
 \node[right]at(6,1.4){图示能量大小；实际感应势取负};
}{线性感应随真实电荷逐渐建立，装配功是三角形面积而非矩形面积。}{ch20:quality:image-energy-a}
图~\ref{ch20:quality:image-energy-a}把真实电荷从零逐渐增加，线性边界值问题使感应势也正比于 $\lambda$。积分 $\int_0^1\lambda\,\dd\lambda=1/2$ 产生感应能中的一半因子；像电荷位置虽可用于求场，却不能被当作独立真实电荷再完整计算一次两体装配能。
}
\newcommand{\ThQPlaneEnergy}{\ThermalFigure{
 \draw[very thick](0,-1.6)--(0,1.6);\foreach\y in {-1.4,-1,...,1.4}{\draw[TextGray](0,\y)--(-.2,\y-.2);}
 \fill[DeepBlue](2,0)circle(3pt)node[above]{$q$};\draw[dashed](4,0)circle(.1)node[above]{$q$（新位置）};
 \draw[<->](0,-.7)--node[below]{$a$}(2,-.7);\draw[<->](0,-1.4)--node[below]{$2a$}(4,-1.4);\node[right]at(5,0){接地无限平面};
}{同一真实电荷由距板 $a$ 移至 $2a$，比较准静态与理想化瞬移两种过程。}{ch20:quality:plane-q}}
\newcommand{\ThUniformSphereBase}[1]{
 \draw[thick](0,0)circle(1.4);\node at(0,-.8){#1};\draw[->](-3.3,-1.5)--(-3.3,1.8)node[above]{$z$};
 \foreach\x in {-2.5,2.5}{\draw[->,DeepBlue,thick](\x,-1)--(\x,1.5);}\node[right]at(3,1){$\vect E_0=E_0\hat{\vect z}$};
 \draw[->](0,0)--(45:1.4)node[right]{$R$};}
\newcommand{\ThQConductingSphere}{\ThermalFigure{\ThUniformSphereBase{中性导体}}{远处原本均匀的外场与导体球；图中箭头只规定外加场方向。}{ch20:quality:conductor-sphere-q}}
\newcommand{\ThADielectricSphere}{\ThermalFigure{\ThUniformSphereBase{$\varepsilon_r$}
 \node[DeepBlue]at(0,1.65){$+$};\node[DeepBlue]at(0,-1.65){$-$};
 \draw[->,thick](.45,-.5)--(.45,.6)node[right]{$\vect E_{\rm in}$};
}{均匀极化球的束缚表面电荷等效产生反向退极化场，球内仍有非零均匀场。}{ch20:quality:dielectric-a}
图~\ref{ch20:quality:dielectric-a}中的正负符号表示束缚电荷的两极，外场与退极化场叠加确定球内场。边界上须连续的是切向 $E$ 和法向 $D$（无自由面电荷），所以不能把导体的 $E_{\rm in}=0$ 边界直接套给有限 $\varepsilon_r$ 的电介质。
}
\newcommand{\ThQDielectricSphere}{\ThermalFigure{\ThUniformSphereBase{电介质 $\varepsilon_r$}}{球外是真空，球内为均匀线性各向同性介质；远场边界为给定匀强场。}{ch20:quality:dielectric-q}}
\newcommand{\ThACapacitors}{\ThermalFigure{
 \draw[very thick](0,0)--(0,2) (1.4,0)--(1.4,2);\draw[->,DeepBlue](.2,1)--(1.2,1)node[above]{$E$};\draw[dashed](-.3,.6)rectangle(.8,1.4);\node[below]at(.7,-.5){平板，薄柱高斯面};
 \begin{scope}[shift={(5,1)}]\draw[thick](0,0)circle(.7);\draw[thick](0,0)circle(1.4);\draw[DeepBlue,dashed](0,0)circle(1.05);\node[below]at(0,-1.7){球形，$a<r<b$};\end{scope}
 \begin{scope}[shift={(10,1)}]\draw[thick](0,0)ellipse(.7 and 1.2);\draw[thick](0,0)ellipse(1.4 and 1.5);\draw[DeepBlue,dashed](0,0)ellipse(1.05 and 1.35);\node[below]at(0,-1.7){柱形，同轴高斯面};\end{scope}
}{先用对称高斯面求间隙电场，再沿场方向积分求极板电势差。}{ch20:quality:capacitors-a}
图~\ref{ch20:quality:capacitors-a}的虚线分别给出三类适用的高斯面。球形通量面积为 $4\pi r^2$，柱形侧面积为 $2\pi rh$，所以其径向电场依次按 $r^{-2}$、$r^{-1}$ 变化；积分 $U=\int_a^bE(r)\dd r$ 后再用 $C=Q/U$，不能直接拿某一点的场强乘所有几何间距。
}
\newcommand{\ThACylinders}{\ThermalFigure{
 \draw[thick](-2.1,0)circle(1);\draw[thick](2.1,0)circle(1);\draw[->](-4,0)--(4,0)node[right]{$x$};
 \fill[DeepBlue](-1.847,0)circle(3pt)node[above]{$-\lambda$};\fill[DeepBlue](1.847,0)circle(3pt)node[above]{$+\lambda$};
 \draw[<->](-2.1,-1.5)--node[below]{$L$}(2.1,-1.5);\node[above]at(0,1.8){等效线电荷在 $x=\pm a$，$a=\sqrt{(L/2)^2-R^2}$};
}{真实圆柱表面是等势圆，辅助线电荷偏离各圆柱几何中心。}{ch20:quality:cylinders-a}
图~\ref{ch20:quality:cylinders-a}中的辅助线电荷放在 $\pm a$，不是 $\pm L/2$。选择这一偏移量才使两圆柱表面分别为固定距离比的等势圆；对数电势差由两个导体的电势相减得到，电容的电荷量按单位轴向长度计算。
}
\newcommand{\ThCapSymbol}[3]{\draw[thick](#1,#2)--++(.75,0) ++(0,-.3)--++(0,.6) ++(.25,0)--++(0,-.6) ++(0,.3)--++(.75,0);\node[above]at(#1+.875,#2+.4){#3};}
\newcommand{\ThAModifiedCapacitor}{\ThermalFigure{
 \draw[thick](0,-1)--(0,1)--(1,1) (0,-1)--(1,-1) (2.75,-1)--(4,-1)--(4,1)--(2.75,1);
 \ThCapSymbol{1}{1}{$C_0/2$}\ThCapSymbol{1}{-1}{$\varepsilon_rC_0/2$}
 \node[below]at(2,-1.8){面积分半：共同电压，并联};
 \begin{scope}[shift={(6,0)}]\ThCapSymbol{0}{0}{$2C_0$}\draw[thick](1.75,0)--(2,0);\ThCapSymbol{2}{0}{$2\varepsilon_rC_0$}\node[below]at(2,-1.8){厚度分半：共同电荷，串联};\end{scope}
}{用与真实几何约束对应的等效网络计算，不按图上左右位置机械判断串并联。}{ch20:quality:modified-a}
图~\ref{ch20:quality:modified-a}左图两部分跨接同一对极板，所以端电压相同；右图两层沿场方向叠放，中间无自由面电荷，故法向 $D$ 相同、对应串联。每层的厚度或面积都要先改写，才得到图中的分电容。
}
\newcommand{\ThAInsertion}{\ThermalFigure{
 \draw[very thick](0,1)--(6,1) (0,-1)--(6,-1);\fill[DeepBlue!18](0,-.9)rectangle(2.7,.9);
 \draw[<->](0,-1.6)--node[below]{$x$}(2.7,-1.6);\draw[->,thick](2.7,0)--(4,0)node[right]{$\dd x>0$};
 \node[right,align=left]at(7,0){定 $Q$：$\delta W_{\rm mech}=-\dd[Q^2/(2C)]$\\定 $U$：电源另供能 $U\dd Q$};
}{介质沿极板方向插入使有效并联面积增大；虚位移与机械功取相同方向。}{ch20:quality:insertion-a}
图~\ref{ch20:quality:insertion-a}把插入方向取为正，因此 $C'(x)>0$。定电荷时由场能下降提供机械功；定电压时还必须计入电源功，机械功为 $U\dd Q-\dd(CU^2/2)$。这两条能量账得到同向的吸入力，不能只对定电压场能取负梯度。
}
\newcommand{\ThPolarizationBase}[1]{
 \draw[thick](0,0)rectangle(3,2);\node at(1.5,1){真空};\draw[thick](5,0)rectangle(8,2);\node at(6.5,1){电介质};
 \foreach\x in {1.5,6.5}{\draw[->,DeepBlue,thick](\x,-.8)--(\x,.1)node[below left]{$E$};}#1}
\newcommand{\ThQPolarizationEnergy}{\ThermalFigure{\ThPolarizationBase{}}{比较相同实际场强下两种介质的能量密度，不能先假定电荷或电压也相同。}{ch20:quality:polarization-q}}
\newcommand{\ThAPolarizationEnergy}{\ThermalFigure{
 \fill[DeepBlue](0,0)circle(3pt)node[left]{$-$};\fill[DeepBlue](2,0)circle(3pt)node[right]{$+$};
 \draw[decorate,decoration={coil,segment length=2mm,amplitude=1.5mm},thick](.2,0)--(1.8,0);
 \draw[->,thick](0,.7)--(2,.7)node[right]{$\vect E,\ \vect p$};\node[below]at(1,-.8){束缚电荷相对位移的示意};
 \node[right,align=left]at(4,0){线性响应：$\vect P\propto\vect E$\\建立极化需逐渐增大外场};
}{微观弹性约束的偶极模型提供极化储能图像；图不表示自由电荷的宏观导电。}{ch20:quality:polarization-a}
图~\ref{ch20:quality:polarization-a}中正负束缚电荷中心被外场拉开，线性响应使建立极化的功含 $1/2$ 因子。比较 $u-u_0=\tfrac12\vect E\cdot\vect P$ 时固定的是同一实际场强；若改成孤立电容器的定 $Q$ 插入过程，场强会改变，能量判断也须重算。
}
\newcommand{\ThQCharging}{\ThermalFigure{
 \draw[thick](0,0)--(0,.8) (0,1.2)--(0,2)--(1,2);\draw[thick](1,1.6)rectangle(2.5,2.4);\node[above]at(1.75,2.6){$R$};
 \draw[thick](2.5,2)--(5,2)--(5,.8) (5,.2)--(5,0)--(0,0);
 \draw[line width=1.2pt](4.5,.8)--(5.5,.8) (4.5,.2)--(5.5,.2);\node[right]at(5.6,.5){$C$};
 \draw[line width=1.2pt](-.4,1.2)--(.4,1.2) (-.25,.8)--(.25,.8);\node[left]at(-.5,1){$\mathcal E$};
}{初始未充电电容经电阻接理想电源；严格零阻时还要检查被忽略的电感等条件。}{ch20:quality:charging-q}}
\newcommand{\ThACharging}{\ThermalFigure{
 \draw[->](0,0)--(6,0)node[right]{$q$};\draw[->](0,0)--(0,3.4)node[above]{每单位电荷的功};
 \draw[thick](0,2.5)--(5,2.5)node[right]{$\mathcal E$};\draw[DeepBlue,thick](0,0)--(5,2.5)node[above left]{$q/C$};
 \fill[DeepBlue!12](0,0)--(5,0)--(5,2.5)--cycle;\node[below]at(5,0){$C\mathcal E$};
}{电源功是矩形面积，电容储能是下方三角形面积，二者之差由电阻耗散。}{ch20:quality:charging-a}
图~\ref{ch20:quality:charging-a}中电源每搬一单位电荷始终做功 $\mathcal E$，而电容反电压随 $q$ 线性增加。两面积的差在任意 $R>0$ 的最终平衡态都相同；减小电阻只缩短过程，不能据此消除这部分积分耗散。
}
\newcommand{\ThAFloatingNode}{\ThermalFigure{
 \coordinate(O)at(0,0);\fill(O)circle(3pt)node[below]{$V$};
 \foreach\a/\i in {90/1,210/2,330/3}{\begin{scope}[rotate=\a]\draw[thick](0,0)--(1,0) (1.25,0)--(2.6,0);\draw[thick](1,-.3)--(1,.3) (1.25,-.3)--(1.25,.3);\end{scope}\node at(\a:3.1){$\varphi_\i,C_\i$};}
 \node[right,align=left]at(4,0){中心三内极板是同一导体\\其净电荷保持初始的零};
}{对浮置中心节点列电荷守恒，先求节点电势，再求三外极板电荷。}{ch20:quality:node-a}
图~\ref{ch20:quality:node-a}中每个内极板电荷为 $C_i(V-\varphi_i)$。三个内极板通过导线相连但未接电源，其电荷和才守恒为零，各内极板的电荷不必分别为零；这给出唯一未知节点电势的线性方程。
}
\newcommand{\ThNetworksBase}[1]{
 \ThCapSymbol{0}{0}{$C_1$}\draw[thick](1.75,0)--(2,0);\ThCapSymbol{2}{0}{$C_2$}\node[left]at(0,0){$A$};\node[right]at(3.75,0){$B$};\node[below]at(2,-.6){串联};
 \begin{scope}[shift={(6,0)}]\draw[thick](0,-1)--(0,1)--(.4,1) (0,-1)--(.4,-1) (2.15,1)--(3,1)--(3,-1)--(2.15,-1);\ThCapSymbol{.4}{1}{$C_1$}\ThCapSymbol{.4}{-1}{$C_2$}\node[left]at(0,0){$A$};\node[right]at(3,0){$B$};\node[below]at(1.5,-1.6){并联};\end{scope}#1}
\newcommand{\ThQCapNetworks}{\ThermalFigure{\ThNetworksBase{}}{比较同一对外端口的两种连接方式，串联中间节点起初中性。}{ch20:quality:networks-q}}
\newcommand{\ThACapNetworks}{\ThermalFigure{\ThNetworksBase{
 \fill(1.88,0)circle(2pt)node[above]{$N$};\node[above]at(5,2.5){串联：节点 $N$ 净电荷零；并联：端口电势共同};
}}{由导线连接关系决定共同电压或共同电荷，再建立两端等效关系。}{ch20:quality:networks-a}
图~\ref{ch20:quality:networks-a}左图的浮置节点 $N$ 只连接两块内极板，故这两块电荷大小相等、符号相反，从而两电容的电荷量大小相等；右图各电容直接接同一对 $A,B$ 端口，所以电压相等。应先证明这两种约束，再分别相加电压或电荷。
}
\newcommand{\ThEarthBase}[1]{
 \fill[DeepBlue!10](0,0)circle(1.5);\draw[thick](0,0)circle(1.5);\draw[thick](0,0)circle(2.1);\node at(0,0){地球 $-Q$};
 \draw[->](0,0)--(150:1.5)node[midway,above]{$R_\oplus$};\draw[<->](1.5,0)--node[above]{$d$}(2.1,0);\node[right]at(3,1){等效外球面 $+Q$};#1}
\newcommand{\ThQEarthCapacitor}{\ThermalFigure{\ThEarthBase{\node[below]at(0,-2.6){大气为有电阻率的球形间隙；厚度示意放大};}}{两同心球面构成电容，大气同时提供漏电通路。}{ch20:quality:earth-q}}
\newcommand{\ThAEarthCapacitor}{\ThermalFigure{\ThEarthBase{
 \draw[DeepBlue,dashed](0,0)circle(1.8);\draw[->,DeepBlue,thick](1.75,.4)--(1.3,.3);
 \node[right,align=left]at(3,-.5){半径 $r$ 处径向通流面积 $4\pi r^2$\\$\dd R=\rho\dd r/(4\pi r^2)$};
}}{同一球形间隙的电容与径向电阻使用相同的几何边界。}{ch20:quality:earth-a}
图~\ref{ch20:quality:earth-a}的虚线球面既可用于高斯定理，也可作为径向漏电流的通流截面。场强与电阻微元都含 $r^{-2}$，因此电容的几何因子恰与电阻的几何因子互为倒数，乘积 $RC=\rho\varepsilon_0$ 与球壳半径无关；雷电补充电荷率应再与漏电流相平衡。
}
